\section{\vrev{시뮬레이션 설계}}
\label{sec:exp_design}
%% ===================================================================

%% -------------------------------------------------------------------
\subsection{\vrev{몬테카를로 시뮬레이션 구성}}
\label{subsec:mc_config}
%% -------------------------------------------------------------------

본 연구의 시뮬레이션은 조건당 $n = 10{,}000$개의 독립 표본을 생성하는
몬테카를로(Monte Carlo) 방법~\citep{rubinstein2016simulation}에 기반한다.
시뮬레이션의 핵심 매개변수는 다음과 같이
기본 복구시간 $R_0 = 6.0$h, 최대 복구시간 $R_{\max} = 24.0$h($= 4R_0$, MTPD 기반 상한),
시그모이드 곡률 $\kappa = 1.2$이다.
각 표본에 대해 입력 변수를 독립적으로 생성하고,
식~(2.10)에 따라 $\DRTO$ 값과
효율성 비율 $\eta = \DRTO / R_0$를 산출한다.

\tabref{tab:sim_parameters}\refpo{tab:sim_parameters}{은}{는} 시뮬레이션의 전체 매개변수를 요약한 것이며, 각 모델 변수의 전체 정의는 2.7.3(변수 및 매개변수 정의)의 [표 2-11]에 정리되어 있다.

\begin{table}[htbp]
  \centering
  \caption{시뮬레이션 매개변수 요약}
  \label{tab:sim_parameters}
  \small
  \begin{tabularx}{\textwidth}{l c c X}
    \toprule
    \textbf{매개변수} & \textbf{기호} & \textbf{값/범위} & \textbf{설정 근거} \\
    \midrule
    기준선 RTO & $R_0$ & 6.0\,h & \nrev{ISO 22301:2019} BIA 기반 표준 설정값 \\
    DRTO 상한 & $R_{\max}$ & 24.0\,h & MTPD 기반 상한 ($= 4R_0$) \\
    곡률 매개변수 & $\kappa$ & 1.2 & 다기준 복합 점수 최대화 기준 최적값 \\
    표본 크기 & $n$ & 10{,}000 & 조건당, CVaR$_{99}$ 안정 추정 기준 \\
    관측성 원시 점수 & $O_{\text{raw}}$ & $\mathrm{Uniform}(0,1)$ & 연속 균등 분포 \\
    관측성 가중치 & $k_O$ & $\mathrm{Uniform}(0,1)$ & 연속 균등 분포 \\
    불확실성 (C2) & $U_{\text{raw}}^{(G)}$ & $N(3,\;1)$ & 경꼬리, $E[U] = 3.0$ \\
    불확실성 (C3) & $U_{\text{raw}}^{(P)}$ & $6 \cdot \mathrm{Pareto}(\alpha\!=\!3)$ & 파레토(중꼬리), $E[U] = 3.0$ \\
    \bottomrule
  \end{tabularx}
\end{table}


\chg{$n = 10{,}000$의 표본 크기는} 다음의 세 가지 기준에 의해 결정되었다.
첫째, 평균 $\bar{\eta}$의 표준오차가 $\mathrm{SE} = \mathrm{SD}/\sqrt{n}$으로서,
C3의 $\eta$ 표준편차 $\mathrm{SD}(\eta_{\text{C3}}) = 0.791$에서도 $\mathrm{SE} = 0.008$에 불과하여
소수점 둘째 자리까지의 정밀도가 보장된다.
둘째, CVaR$_{99}$의 산출에 상위 $1\%$에 해당하는 100개 이상의 표본이
필요하며, $n = 10{,}000$은 이 요구를 충족한다.
셋째, 부분 표본 분석에서 $n = 5{,}000$ 이상일 때 핵심 통계량이
소수점 셋째 자리까지 수렴함을 사전 확인하였으며 4.6(수렴 진단 및 재현성)에서 그 상세 결과를 알 수 있다.


%% -------------------------------------------------------------------
\subsection{\vrev{난수 생성 방법}}
\label{subsec:rng_method}
%% -------------------------------------------------------------------

본 시뮬레이션의 난수 생성은 [표 3-3]의 시뮬레이션 매개변수 요약 기준에 따라 설계되었으며, 난수 생성기의 상세 규격은 [표 3-4]에 정리하였다.

\begin{table}[htbp]
  \centering
  \caption{난수 생성 방법 규격}
  \label{tab:rng_spec}
  \small
  \begin{tabularx}{\textwidth}{l X}
    \toprule
    \textbf{항목} & \textbf{규격 및 근거} \\
    \midrule
    난수 생성기 & PCG64 (Permuted Congruential Generator, 64-bit). NumPy 1.17+의 \texttt{default\_rng()} 사용.
      기존 Mersenne Twister(MT19937) 대비 통계적 품질·속도·시드 공간 우수. \\
    \m{주기(Period)} & $2^{128}$ (MT19937의 $2^{19937}$보다 작으나,
      $n = 10{,}000$ 수준의 시뮬레이션에서 주기 소진 위험 없음) \\
    시드 지정 & 변수별 개별 시드를 명시적으로 지정하여 \texttt{default\_rng(seed)} 호출.
      암묵적 시스템 시각 기반 시드를 사용하지 않음. (명시 미지정 시 시스템 시각 기반 자동 할당으로 재현 불가) \\
    재현성 보장 & 동일 시드$\cdot$동일 알고리즘$\cdot$동일 호출 순서에서 비트 수준 동일 결과 보장.
      NumPy 버전 간 호환성을 위해 \texttt{Generator} API \m{사용(legacy \texttt{RandomState} 미사용)} \\
    분포 생성 & 균등분포: \texttt{rng.uniform(0, 1)},
      정규분포: \texttt{rng.normal(3, 1)},
      파레토분포: \tnew{\texttt{6 * rng.pareto(3)}} ($E[U] = 3.0$ 보장) \\
    \bottomrule
  \end{tabularx}
\end{table}

\chg{여기서 세 가지 설계 근거는 다음과 같다. 시드 지정과 관련하여, 시드를 명시하지 않으면 많은 난수 생성기가 현재 시스템 시각(system clock)을 시드로 자동 할당하여 실행 시점마다 결과가 달라지므로, 본 연구는 모든 확률 스트림에 고정 시드를 명시적으로 지정하였다. 재현성과 관련하여, NumPy의 구 API인 \texttt{RandomState}는 Mersenne Twister 엔진을 사용하여 버전 갱신 시 동일 시드에서도 난수열이 달라질 수 있는 반면, 신 API인 \texttt{Generator}(NumPy 1.17 이상)는 버전 간 비트 수준 재현성을 보장하므로 이를 채택하였다. 분포 생성과 관련하여, NumPy의 \texttt{Generator.pareto(a)}는 Lomax(파레토 II형) 분포를 반환하여 $E[\mathrm{Lomax}(\alpha=3)] = 1/(\alpha-1) = 0.5$이므로 6을 곱한 \texttt{6 * rng.pareto(3)}의 기댓값이 $3.0$이 된다.}


%% -------------------------------------------------------------------
\subsection{\vrev{비중첩 대역 시드 체계}}
\label{subsec:seed_system}
%% -------------------------------------------------------------------

시뮬레이션의 완전한 재현성을 보장하면서 변수 간 통계적 독립성을
유지하기 위해, 마스터 시드 $s = 42$에서 변수별 시드를 비중첩 대역(non-overlapping band)으로
파생하는 체계를 설계하였다. \tabref{tab:seed_system}\refpo{tab:seed_system}{은}{는}
각 변수의 시드 값과 대역 범위를 정리한 것이다.

\begin{table}[htbp]
  \centering
  \caption{비중첩 대역 시드 \m{체계(마스터 시드 $s = 42$)}}
  \label{tab:seed_system}
  \begin{tabular}{llcc}
    \toprule
    변수 & 기호 & 시드 값 & 대역 범위 \\
    \midrule
    관측성 원시 점수 & $O_{\text{raw}}$ & 42 & $[0\text{--}9999]$ \\
    관측성 가중치 & $k_O$ & 20{,}042 & $[20000\text{--}29999]$ \\
    \m{불확실성(C2, Gaussian)} & $U^{(G)}$ & 30{,}042 & $[30000\text{--}39999]$ \\
    \m{불확실성(C3, Pareto)} & $U^{(P)}$ & 40{,}042 & $[40000\text{--}49999]$ \\
    \bottomrule
  \end{tabular}
\end{table}

각 대역은 $10{,}000$개의 고유 시드를 수용하며, 대역 간 간격이 충분하여
난수 생성기(NumPy PCG64)의 상태 간 교차 오염을 구조적으로
방지한다. 마스터 시드를 변경하면 모든 변수의 시드가 동일한 오프셋 규칙으로
갱신되므로, 다중 시드 검증($s = 0, 1, \ldots, 9{,}999$) 시에도
일관된 독립성이 유지된다.
본 시드 체계의 독립성은 \pdferrR{표~\ref{tab:independence}}에서 실증적으로 검증되며 H6을 충족한다.


%% -------------------------------------------------------------------
\subsection{\padd{재현성 자료 공개}}
\label{subsec:reproducibility-release}
%% -------------------------------------------------------------------

\padd{본 연구의 후속 \chg{검증과 확장} 가능성을 확보하기 위해, 시뮬레이션 코드와
데이터를 GitHub 공개 저장소에 공개한다. \chg{저장소의 공개 주소는 \texttt{https://github.com/WindoorsLEE/drto-dynamic-rto}이며, 코드는 MIT 라이선스로, 데이터는 CC BY 4.0 라이선스로 배포한다.} \chg{공개 자료는 네 종류로 구성된다. 첫째, 시뮬레이션 코드(Python 3.12와 NumPy 2.4 기반)는 C0에서 C3까지의 조건별 몬테카를로 생성과 분할 회귀 및 CVaR 산출, 그리고 $10^8$회 다중 시드 강건성 검증 스크립트를 포함한다. 둘째, 원시 난수열(\texttt{raw\_data/})은 비중첩 대역 시드 체계로 생성된 $O_{\text{raw}}$, $k_O$, $U^{(G)}$, $U^{(P)}$ 4종을 표본 크기 $n = 10{,}000$의 CSV로 담는다. 셋째, 계산 결과(\texttt{calc\_data/})는 조건별 $\DRTO$ 값과 평균, 표준편차, 백분위 등 기준값을 정리한 \texttt{statistics.json}으로 이루어진다. 넷째, LaTeX 소스는 본 논문 본문과 부록을 빌드 가능한 형태로 제공한다.}
재현 절차(\texttt{README.md} 참조)는 저장소 클론 후 \texttt{make all} 실행 시
본 연구의 모든 수치 결과(조건별 \chg{평균, CDF, 회귀계수} 등)가 비트 수준에서
동일하게 재현된다.}


%% -------------------------------------------------------------------
\subsection{\vrev{변수 간 독립성 검증}}
\label{subsec:independence}
%% -------------------------------------------------------------------

비중첩 대역 시드 체계의 유효성을 확인하기 위해, 생성된 $n = 10{,}000$개
표본에서 변수 쌍 간 피어슨 상관계수를 산출하였다.
\tabref{tab:independence}에 제시된 바와 같이, 모든 변수 쌍의
상관계수 절댓값이 $|r| < 0.02$로 나타나 실질적 독립성이 확인되었다. \chg{여기서 기준값 $0.02$는 표본 크기 $n = 10{,}000$에서 귀무가설 $H_0\!: \rho = 0$ 하의 상관계수 표준오차가 $\mathrm{SE}(r) \approx 1/\sqrt{n} = 0.01$이고 95\% 신뢰구간이 $\pm 1.96 \times 0.01 \approx \pm 0.020$이라는 데 근거하며, 따라서 $|r| < 0.02$는 관측된 상관이 유의수준 $\alpha = 0.05$에서 영과 통계적으로 구별되지 않음을 의미한다.}

\begin{table}[htbp]
  \centering
  \caption{입력 변수 간 피어슨 상관계수 ($n = 10{,}000$)}
  \label{tab:independence}
  \begin{tabular}{lcc}
    \toprule
    변수 쌍 & 상관계수 $r$ & 독립성 판정 \\
    \midrule
    $k_O$ -- $O_{\text{raw}}$ & $0.016$ & $|r| < 0.02$ \\
    $O_{\text{raw}}$ -- $U^{(G)}$ & $0.009$ & $|r| < 0.02$ \\
    $O_{\text{raw}}$ -- $U^{(P)}$ & $-0.001$ & $|r| < 0.02$ \\
    \bottomrule
  \end{tabular}
\end{table}


%% -------------------------------------------------------------------
\subsection{\vrev{분포 설정 및 \imod{공정 비교}{공정한 비교} 설계}}
\label{subsec:dist_settings}
%% -------------------------------------------------------------------

C2 조건과 C3 조건의 불확실성 분포는 평균이 동일하되 꼬리 특성이
상이하도록 설계하였다. 이를 통해 분포 형태(shape)가 $\DRTO$에 미치는
고유한 효과를 분리하여 관측할 수 있다.

\chg{C2 조건의 가우시안 불확실성은} $U_{\text{raw}} \sim N(\mu=3,\;\sigma=1)$에서 추출한다.
시그모이드 입력 $z_U = \kappa \cdot k_U \cdot R_0 \cdot U_{\text{raw}} / (R_{\max} - R_0)$로
변환되며, $E[U_{\text{raw}}] = 3.0$의 기대값을 가진다. 시뮬레이션에서
관측된 통계량은 $\bar{U}_{\text{C2}} = 3.003$, $\mathrm{SD} = 0.996$,
$P_{99} = 5.333$이다. 가우시안 분포의 특성상 극단값이 평균에서
크게 이탈하지 않으며, $P_{99}/\mu$ 비율은 $1.78$에 불과하다. \chg{여기서 $P_{99}/\mu$ 비율은 99번째 백분위수와 모평균의 비율로서 분포의 우측 꼬리가 중심으로부터 얼마나 이탈하는지를 나타내며, 본 시뮬레이션에서 $\mu$는 추정값이 아니라 설계 파라미터 $E[U_{\text{raw}}] = 3.0$이므로 정확한 모집단 평균을 사용한다.}

\chg{C3 조건의 파레토 불확실성은} $U_{\text{raw}} \sim 6 \cdot \mathrm{Pareto}(\alpha=3)$에서 추출한다.
형상모수 $\alpha = 3$은 2차 모멘트가 존재하되($\alpha > 2$),
파레토(Heavy-Tail) 특성을 보유하는 분포를 생성한다.
시뮬레이션에서 관측된 통계량은 $\bar{U}_{\text{C3}} = 3.019$,
$\mathrm{SD} = 5.303$, $P_{99} = 23.402$이다.
$P_{99}/\mu$ 비율은 $23.402 / 3.0 = 7.80$으로서,
\textbf{가우시안의 $1.78$ 대비 $4.39$배}에 달하며,
극단 사건의 발생 빈도와 규모가 현저히 크다는 것을 확인할 수 있다.

두 분포의 평균 $\bar{U}_{\text{raw}} \approx 3.0$이 실질적으로
동일하므로, C2와 C3 간 $\eta$ 차이는 불확실성의 크기가 아닌
분포 형태, \chg{특히 우측 꼬리의 중량에} 기인하는 것으로 해석할 수 있다.

\chg{이상에서 제시한 연구 설계와 시뮬레이션 방법에 따라 수행한 실증 분석의 결과는 다음 장에서 조건별 기술통계, 회귀 및 분할 회귀, 강건성 검증의 순으로 제시한다.}
